\chapter{Répertoire du formulaire}\label{ch:rep_formulaire}

Ce répertoire détaille, élément par élément, les données résumées au tableau~\ref{tab:formulaire}. Pour chaque élément, on donne la règle de repérage (expression régulière appliquée au texte normalisé, \emph{u/v} et \emph{i/j} unifiés, coupures de ligne neutralisées), l’effectif de l’ensemble C avec et sans restitution, la répartition par période des attestations sans restitution, les formes relevées et la liste des notices où l’élément est attesté sans restitution.

Le répertoire est entièrement produit par le programme de génération des tableaux (\texttt{04\_generated.py}) à partir du champ « Texte transmis ». Il est donc reproductible et contrôlable : une erreur de règle se corrige dans le programme, et tous les effectifs du livre sont recalculés. Les règles ont été vérifiées sur les formes relevées ; deux règles trop larges ont été corrigées à cette occasion (celle de \emph{filius} retenait le nom \emph{Filippus}, celle de \emph{coniux} le verbe \emph{comparaverunt}). Les effectifs diffèrent donc de ceux de la version précédente pour ces deux éléments (\emph{filius / filia} : 33 textes avec restitution au lieu de 42). Les effectifs par période se rapportent aux textes de l’ensemble C datés de la période, quel que soit le fondement de l’année ; le chapitre~\ref{ch:mesures} examine la sensibilité des principaux résultats à ce choix.

{\small
% généré par scripts/04_generated.py
\subsection*{in pace}
\addcontentsline{toc}{subsection}{in pace}
\noindent Règle (texte normalisé, \textit{u/v} et \textit{i/j} unifiés) : \texttt{\small \textbackslash{}bin\textbackslash{}s+pace\textbackslash{}b}. Ensemble C : 193 textes avec restitutions, 147 sans restitution ; écart dû aux restitutions : 46.\par
\noindent Par période (attestations sans restitution / textes de la période) : 340–359 : 21/55 (38,2~\%) ; 360–374 : 46/121 (38,0~\%) ; 375–389 : 38/135 (28,1~\%) ; 390–403 : 42/154 (27,3~\%).\par
\noindent Formes relevées (occurrences) : in pace (204), in   pace (4).\par
\noindent{\small Notices (élément attesté sans restitution) : 4, 9, 10, 16, 17, 20, 22, 26, 28, 33, 34, 35, 38, 41, 42, 45, 46, 47, 49, 54, 55, 56, 57, 59, 61, 63, 66, 68, 70, 72, 75, 76, 78, 80, 81, 83, 94, 99, 102, 104, 115, 117, 119, 120, 121, 124, 125, 126, 127, 128, 130, 132, 133, 134, 139, 140, 143, 144, 145, 150, 151, 155, 156, 157, 166, 168, 172, 179, 183, 184, 185, 191, 196, 200, 202, 203, 205, 207, 210, 211, 213, 221, 223, 231, 233, 234, 237, 240, 241, 250, 260, 261, 271, 276, 279, 284, 287, 288, 291, 295, 296, 298, 300, 303, 304, 320, 323, 336, 346, 347, 351, 353, 359, 361, 365, 367, 368, 369, 371, 374, 376, 381, 383, 388, 389, 396, 399, 402, 403, 404, 409, 411, 416, 419, 420, 421, 427, 431, 444, 446, 448, 451, 453, 460, 463, 469, 470.}\par\medskip
\subsection*{depositio / depositus}
\addcontentsline{toc}{subsection}{depositio / depositus}
\noindent Règle (texte normalisé, \textit{u/v} et \textit{i/j} unifiés) : \texttt{\small \textbackslash{}bdepo?s+\textbackslash{}w*}. Ensemble C : 247 textes avec restitutions, 179 sans restitution ; écart dû aux restitutions : 68.\par
\noindent Par période (attestations sans restitution / textes de la période) : 340–359 : 24/55 (43,6~\%) ; 360–374 : 50/121 (41,3~\%) ; 375–389 : 48/135 (35,6~\%) ; 390–403 : 57/154 (37,0~\%).\par
\noindent Formes relevées (occurrences) : depositus (127), deposita (112), depositio (13), depossio (3), deposit (2), depositi (1), depos (1), deposito (1), depositae (1).\par
\noindent{\small Notices (élément attesté sans restitution) : 4, 12, 13, 16, 17, 18, 20, 22, 23, 25, 26, 34, 35, 38, 40, 41, 44, 46, 47, 48, 49, 50, 54, 55, 57, 59, 60, 63, 66, 67, 68, 73, 75, 76, 77, 78, 81, 83, 90, 97, 99, 102, 105, 109, 111, 115, 117, 119, 124, 125, 126, 127, 128, 132, 133, 137, 139, 140, 142, 143, 146, 151, 154, 155, 156, 157, 158, 166, 167, 169, 171, 172, 175, 179, 181, 183, 185, 188, 191, 195, 196, 199, 200, 202, 204, 205, 206, 207, 208, 210, 211, 213, 225, 226, 231, 233, 235, 241, 242, 243, 244, 251, 255, 258, 261, 265, 272, 279, 286, 288, 291, 292, 295, 297, 298, 301, 302, 303, 304, 310, 311, 318, 319, 320, 323, 324, 325, 336, 345, 346, 347, 351, 353, 355, 357, 364, 365, 366, 367, 368, 370, 371, 372, 373, 376, 378, 383, 385, 386, 387, 388, 391, 392, 395, 396, 399, 402, 404, 405, 409, 411, 412, 413, 418, 420, 421, 434, 441, 443, 446, 448, 449, 451, 459, 460, 463, 466, 469, 470.}\par\medskip
\subsection*{vixit (et variantes)}
\addcontentsline{toc}{subsection}{vixit (et variantes)}
\noindent Règle (texte normalisé, \textit{u/v} et \textit{i/j} unifiés) : \texttt{\small \textbackslash{}b(\allowbreak{}uix\textbackslash{}w*|\allowbreak{}bix\textbackslash{}w*|\allowbreak{}uiuit|\allowbreak{}uibit)\textbackslash{}b}. Ensemble C : 264 textes avec restitutions, 189 sans restitution ; écart dû aux restitutions : 75.\par
\noindent Par période (attestations sans restitution / textes de la période) : 340–359 : 26/55 (47,3~\%) ; 360–374 : 53/121 (43,8~\%) ; 375–389 : 53/135 (39,3~\%) ; 390–403 : 57/154 (37,0~\%).\par
\noindent Formes relevées (occurrences) : uixit (277), uiuit (4), uix (1).\par
\noindent{\small Notices (élément attesté sans restitution) : 2, 4, 6, 7, 9, 13, 15, 16, 17, 19, 20, 22, 25, 26, 33, 34, 35, 38, 40, 41, 42, 45, 47, 49, 53, 54, 57, 59, 60, 61, 63, 68, 70, 72, 75, 79, 81, 83, 84, 90, 91, 94, 95, 99, 100, 102, 109, 111, 115, 117, 118, 119, 122, 124, 125, 126, 128, 130, 131, 132, 133, 136, 140, 142, 143, 144, 145, 146, 147, 149, 151, 156, 157, 163, 166, 168, 169, 170, 172, 181, 183, 184, 185, 186, 188, 189, 191, 193, 197, 200, 201, 202, 203, 207, 208, 210, 211, 215, 219, 222, 223, 225, 226, 231, 235, 237, 244, 245, 251, 255, 260, 263, 268, 269, 271, 272, 275, 276, 278, 288, 292, 296, 297, 298, 299, 301, 303, 304, 306, 308, 310, 317, 322, 324, 334, 336, 338, 339, 343, 346, 347, 349, 351, 353, 354, 355, 356, 363, 365, 366, 367, 368, 372, 374, 376, 382, 383, 385, 388, 390, 391, 396, 398, 399, 400, 401, 403, 405, 409, 412, 418, 419, 427, 429, 431, 433, 440, 442, 443, 446, 448, 450, 463, 466, 467, 468, 469, 472, 474.}\par\medskip
\subsection*{bene merenti}
\addcontentsline{toc}{subsection}{bene merenti}
\noindent Règle (texte normalisé, \textit{u/v} et \textit{i/j} unifiés) : \texttt{\small \textbackslash{}bbene\textbackslash{}s*mere\textbackslash{}w*}. Ensemble C : 86 textes avec restitutions, 53 sans restitution ; écart dû aux restitutions : 33.\par
\noindent Par période (attestations sans restitution / textes de la période) : 340–359 : 7/55 (12,7~\%) ; 360–374 : 23/121 (19,0~\%) ; 375–389 : 10/135 (7,4~\%) ; 390–403 : 13/154 (8,4~\%).\par
\noindent Formes relevées (occurrences) : bene merenti (84), bene   merenti (5), bene merente (1).\par
\noindent{\small Notices (élément attesté sans restitution) : 4, 16, 17, 33, 35, 38, 47, 70, 75, 76, 77, 99, 100, 109, 115, 117, 124, 127, 128, 131, 132, 139, 140, 143, 144, 145, 152, 155, 166, 172, 198, 210, 223, 226, 260, 271, 276, 295, 300, 304, 319, 320, 324, 336, 351, 365, 376, 396, 399, 409, 428, 431, 444.}\par\medskip
\subsection*{fecit / fecerunt}
\addcontentsline{toc}{subsection}{fecit / fecerunt}
\noindent Règle (texte normalisé, \textit{u/v} et \textit{i/j} unifiés) : \texttt{\small \textbackslash{}bfec(\allowbreak{}it|\allowbreak{}erunt|\allowbreak{}i)\textbackslash{}b}. Ensemble C : 74 textes avec restitutions, 61 sans restitution ; écart dû aux restitutions : 13.\par
\noindent Par période (attestations sans restitution / textes de la période) : 340–359 : 10/55 (18,2~\%) ; 360–374 : 16/121 (13,2~\%) ; 375–389 : 19/135 (14,1~\%) ; 390–403 : 16/154 (10,4~\%).\par
\noindent Formes relevées (occurrences) : fecit (70), fecerunt (8), feci (1).\par
\noindent{\small Notices (élément attesté sans restitution) : 4, 6, 21, 22, 26, 33, 36, 42, 46, 55, 59, 68, 73, 83, 104, 110, 114, 115, 126, 132, 133, 135, 144, 146, 147, 148, 184, 185, 193, 194, 197, 200, 204, 207, 223, 235, 247, 257, 262, 265, 286, 295, 296, 297, 300, 319, 320, 347, 353, 364, 365, 368, 373, 374, 386, 399, 409, 420, 428, 441, 454.}\par\medskip
\subsection*{hic (iacet, positus, requiescit…)}
\addcontentsline{toc}{subsection}{hic (iacet, positus, requiescit…)}
\noindent Règle (texte normalisé, \textit{u/v} et \textit{i/j} unifiés) : \texttt{\small \textbackslash{}bhic\textbackslash{}b}. Ensemble C : 33 textes avec restitutions, 27 sans restitution ; écart dû aux restitutions : 6.\par
\noindent Par période (attestations sans restitution / textes de la période) : 340–359 : 1/55 (1,8~\%) ; 360–374 : 4/121 (3,3~\%) ; 375–389 : 9/135 (6,7~\%) ; 390–403 : 13/154 (8,4~\%).\par
\noindent Formes relevées (occurrences) : hic (34).\par
\noindent{\small Notices (élément attesté sans restitution) : 22, 60, 78, 95, 151, 213, 214, 244, 255, 269, 298, 299, 303, 318, 333, 352, 374, 388, 390, 396, 400, 428, 433, 450, 460, 466, 470.}\par\medskip
\subsection*{requiescit / quiescit}
\addcontentsline{toc}{subsection}{requiescit / quiescit}
\noindent Règle (texte normalisé, \textit{u/v} et \textit{i/j} unifiés) : \texttt{\small \textbackslash{}b(\allowbreak{}re)?quie(\allowbreak{}sc|\allowbreak{}u)\textbackslash{}w*}. Ensemble C : 37 textes avec restitutions, 29 sans restitution ; écart dû aux restitutions : 8.\par
\noindent Par période (attestations sans restitution / textes de la période) : 340–359 : 2/55 (3,6~\%) ; 360–374 : 8/121 (6,6~\%) ; 375–389 : 5/135 (3,7~\%) ; 390–403 : 14/154 (9,1~\%).\par
\noindent Formes relevées (occurrences) : quiescit (15), requiescit (8), quiescet (3), quiescis (3), requieuit (2), requiescet (2), quiesces (1), requiescunt (1), quiescenti (1), quieuit (1), quiescas (1).\par
\noindent{\small Notices (élément attesté sans restitution) : 22, 42, 59, 70, 72, 73, 80, 102, 125, 145, 203, 204, 213, 214, 231, 333, 371, 374, 378, 388, 391, 404, 405, 406, 446, 450, 460, 467, 474.}\par\medskip
\subsection*{decessit / recessit / defunctus}
\addcontentsline{toc}{subsection}{decessit / recessit / defunctus}
\noindent Règle (texte normalisé, \textit{u/v} et \textit{i/j} unifiés) : \texttt{\small \textbackslash{}b(\allowbreak{}de|\allowbreak{}re)cess\textbackslash{}w*|\allowbreak{}\textbackslash{}bdefunct\textbackslash{}w*}. Ensemble C : 36 textes avec restitutions, 28 sans restitution ; écart dû aux restitutions : 8.\par
\noindent Par période (attestations sans restitution / textes de la période) : 340–359 : 7/55 (12,7~\%) ; 360–374 : 8/121 (6,6~\%) ; 375–389 : 9/135 (6,7~\%) ; 390–403 : 4/154 (2,6~\%).\par
\noindent Formes relevées (occurrences) : decessit (17), defunctus (9), recessit (7), defuncta (5).\par
\noindent{\small Notices (élément attesté sans restitution) : 4, 7, 20, 28, 31, 33, 54, 58, 92, 118, 119, 120, 136, 139, 160, 193, 194, 197, 231, 260, 269, 276, 300, 308, 381, 430, 449, 478.}\par\medskip
\subsection*{dormitio / dormit}
\addcontentsline{toc}{subsection}{dormitio / dormit}
\noindent Règle (texte normalisé, \textit{u/v} et \textit{i/j} unifiés) : \texttt{\small \textbackslash{}bdorm\textbackslash{}w*}. Ensemble C : 6 textes avec restitutions, 3 sans restitution ; écart dû aux restitutions : 3.\par
\noindent Par période (attestations sans restitution / textes de la période) : 340–359 : 0/55 (0,0~\%) ; 360–374 : 2/121 (1,7~\%) ; 375–389 : 0/135 (0,0~\%) ; 390–403 : 1/154 (0,6~\%).\par
\noindent Formes relevées (occurrences) : dormitio (3), dormit (3), dormientes (1).\par
\noindent{\small Notices (élément attesté sans restitution) : 73, 134, 400.}\par\medskip
\subsection*{locus}
\addcontentsline{toc}{subsection}{locus}
\noindent Règle (texte normalisé, \textit{u/v} et \textit{i/j} unifiés) : \texttt{\small \textbackslash{}bloc(\allowbreak{}us|\allowbreak{}um|\allowbreak{}o|\allowbreak{}i)\textbackslash{}b}. Ensemble C : 26 textes avec restitutions, 19 sans restitution ; écart dû aux restitutions : 7.\par
\noindent Par période (attestations sans restitution / textes de la période) : 340–359 : 2/55 (3,6~\%) ; 360–374 : 4/121 (3,3~\%) ; 375–389 : 0/135 (0,0~\%) ; 390–403 : 13/154 (8,4~\%).\par
\noindent Formes relevées (occurrences) : locum (15), locus (10), loco (1).\par
\noindent{\small Notices (élément attesté sans restitution) : 20, 48, 73, 130, 148, 150, 320, 347, 365, 368, 379, 387, 395, 396, 400, 411, 454, 464, 478.}\par\medskip
\subsection*{bisomus / trisomus / quadrisomus}
\addcontentsline{toc}{subsection}{bisomus / trisomus / quadrisomus}
\noindent Règle (texte normalisé, \textit{u/v} et \textit{i/j} unifiés) : \texttt{\small \textbackslash{}b(\allowbreak{}bi|\allowbreak{}tri|\allowbreak{}quadri)som\textbackslash{}w*}. Ensemble C : 10 textes avec restitutions, 6 sans restitution ; écart dû aux restitutions : 4.\par
\noindent Par période (attestations sans restitution / textes de la période) : 340–359 : 0/55 (0,0~\%) ; 360–374 : 1/121 (0,8~\%) ; 375–389 : 1/135 (0,7~\%) ; 390–403 : 4/154 (2,6~\%).\par
\noindent Formes relevées (occurrences) : bisomum (6), bisomus (2), quadrisomum (1), trisomus (1).\par
\noindent{\small Notices (élément attesté sans restitution) : 130, 275, 365, 374, 395, 478.}\par\medskip
\subsection*{achat (emere, comparare)}
\addcontentsline{toc}{subsection}{achat (emere, comparare)}
\noindent Règle (texte normalisé, \textit{u/v} et \textit{i/j} unifiés) : \texttt{\small \textbackslash{}b(\allowbreak{}emit|\allowbreak{}emerunt|\allowbreak{}empt\textbackslash{}w*|\allowbreak{}comparau\textbackslash{}w*|\allowbreak{}comparat\textbackslash{}w*|\allowbreak{}comparass\textbackslash{}w*)\textbackslash{}b}. Ensemble C : 7 textes avec restitutions, 4 sans restitution ; écart dû aux restitutions : 3.\par
\noindent Par période (attestations sans restitution / textes de la période) : 340–359 : 0/55 (0,0~\%) ; 360–374 : 0/121 (0,0~\%) ; 375–389 : 0/135 (0,0~\%) ; 390–403 : 4/154 (2,6~\%).\par
\noindent Formes relevées (occurrences) : comparauerunt (2), emerunt (1), emit (1), comparauit (1), emptum (1), emptus (1).\par
\noindent{\small Notices (élément attesté sans restitution) : 400, 406, 411, 478.}\par\medskip
\subsection*{fossor}
\addcontentsline{toc}{subsection}{fossor}
\noindent Règle (texte normalisé, \textit{u/v} et \textit{i/j} unifiés) : \texttt{\small \textbackslash{}bfossor\textbackslash{}w*}. Ensemble C : 1 textes avec restitutions, 1 sans restitution ; écart dû aux restitutions : 0.\par
\noindent Par période (attestations sans restitution / textes de la période) : 340–359 : 0/55 (0,0~\%) ; 360–374 : 0/121 (0,0~\%) ; 375–389 : 0/135 (0,0~\%) ; 390–403 : 1/154 (0,6~\%).\par
\noindent Formes relevées (occurrences) : fossore (1).\par
\noindent{\small Notices (élément attesté sans restitution) : 478.}\par\medskip
\subsection*{coniux / compar / maritus / uxor}
\addcontentsline{toc}{subsection}{coniux / compar / maritus / uxor}
\noindent Règle (texte normalisé, \textit{u/v} et \textit{i/j} unifiés) : \texttt{\small \textbackslash{}b(\allowbreak{}coniu[gnx]\textbackslash{}w*|\allowbreak{}compar(\allowbreak{}i|\allowbreak{}e|\allowbreak{}em)?|\allowbreak{}marit\textbackslash{}w*|\allowbreak{}uxor\textbackslash{}w*)\textbackslash{}b}. Ensemble C : 71 textes avec restitutions, 45 sans restitution ; écart dû aux restitutions : 26.\par
\noindent Par période (attestations sans restitution / textes de la période) : 340–359 : 12/55 (21,8~\%) ; 360–374 : 13/121 (10,7~\%) ; 375–389 : 12/135 (8,9~\%) ; 390–403 : 8/154 (5,2~\%).\par
\noindent Formes relevées (occurrences) : marito (22), coniugi (16), compari (6), uxor (5), maritus (5), coniux (5), compare (4), coniuge (4), uxori (3), uxore (3), coniunx (2), uxoris (1).\par
\noindent{\small Notices (élément attesté sans restitution) : 23, 26, 29, 32, 33, 34, 36, 42, 43, 46, 49, 53, 56, 77, 104, 110, 114, 115, 117, 125, 130, 133, 144, 146, 147, 184, 185, 193, 200, 223, 244, 275, 288, 297, 304, 306, 310, 319, 320, 324, 365, 382, 386, 399, 400.}\par\medskip
\subsection*{virginius / virginia}
\addcontentsline{toc}{subsection}{virginius / virginia}
\noindent Règle (texte normalisé, \textit{u/v} et \textit{i/j} unifiés) : \texttt{\small \textbackslash{}b[bu]irgini(\allowbreak{}us|\allowbreak{}a|\allowbreak{}o|\allowbreak{}um|\allowbreak{}ae)\textbackslash{}b}. Ensemble C : 5 textes avec restitutions, 3 sans restitution ; écart dû aux restitutions : 2.\par
\noindent Par période (attestations sans restitution / textes de la période) : 340–359 : 1/55 (1,8~\%) ; 360–374 : 0/121 (0,0~\%) ; 375–389 : 2/135 (1,5~\%) ; 390–403 : 0/154 (0,0~\%).\par
\noindent Formes relevées (occurrences) : uirginia (3), uirginio (2), uirginium (1).\par
\noindent{\small Notices (élément attesté sans restitution) : 4, 218, 265.}\par\medskip
\subsection*{filius / filia}
\addcontentsline{toc}{subsection}{filius / filia}
\noindent Règle (texte normalisé, \textit{u/v} et \textit{i/j} unifiés) : \texttt{\small \textbackslash{}bfili(\allowbreak{}us|\allowbreak{}a|\allowbreak{}o|\allowbreak{}ae|\allowbreak{}i|\allowbreak{}is|\allowbreak{}um|\allowbreak{}am|\allowbreak{}e)\textbackslash{}b}. Ensemble C : 33 textes avec restitutions, 27 sans restitution ; écart dû aux restitutions : 6.\par
\noindent Par période (attestations sans restitution / textes de la période) : 340–359 : 3/55 (5,5~\%) ; 360–374 : 9/121 (7,4~\%) ; 375–389 : 9/135 (6,7~\%) ; 390–403 : 6/154 (3,9~\%).\par
\noindent Formes relevées (occurrences) : filio (15), filiae (12), filia (5), filius (2), filiis (1).\par
\noindent{\small Notices (élément attesté sans restitution) : 13, 35, 42, 59, 68, 81, 84, 92, 99, 100, 127, 143, 186, 196, 211, 213, 233, 259, 268, 297, 306, 347, 353, 378, 399, 412, 431.}\par\medskip
\subsection*{parentes / pater / mater}
\addcontentsline{toc}{subsection}{parentes / pater / mater}
\noindent Règle (texte normalisé, \textit{u/v} et \textit{i/j} unifiés) : \texttt{\small \textbackslash{}b(\allowbreak{}parent\textbackslash{}w*|\allowbreak{}pater|\allowbreak{}patri|\allowbreak{}patre|\allowbreak{}patris|\allowbreak{}patrem|\allowbreak{}mater|\allowbreak{}matri|\allowbreak{}matre|\allowbreak{}matris|\allowbreak{}matrem)\textbackslash{}b}. Ensemble C : 29 textes avec restitutions, 24 sans restitution ; écart dû aux restitutions : 5.\par
\noindent Par période (attestations sans restitution / textes de la période) : 340–359 : 4/55 (7,3~\%) ; 360–374 : 7/121 (5,8~\%) ; 375–389 : 7/135 (5,2~\%) ; 390–403 : 6/154 (3,9~\%).\par
\noindent Formes relevées (occurrences) : parentes (9), pater (6), mater (4), matri (3), parentibus (2), parentis (1), patri (1), parentum (1), patris (1), matre (1), patre (1), parenti (1).\par
\noindent{\small Notices (élément attesté sans restitution) : 2, 17, 22, 35, 64, 68, 73, 81, 102, 124, 126, 203, 241, 244, 268, 296, 297, 306, 353, 382, 409, 439, 461, 470.}\par\medskip
\subsection*{dulcissimus / carissimus}
\addcontentsline{toc}{subsection}{dulcissimus / carissimus}
\noindent Règle (texte normalisé, \textit{u/v} et \textit{i/j} unifiés) : \texttt{\small \textbackslash{}b(\allowbreak{}dulciss\textbackslash{}w*|\allowbreak{}cariss\textbackslash{}w*|\allowbreak{}kariss\textbackslash{}w*)}. Ensemble C : 30 textes avec restitutions, 24 sans restitution ; écart dû aux restitutions : 6.\par
\noindent Par période (attestations sans restitution / textes de la période) : 340–359 : 5/55 (9,1~\%) ; 360–374 : 10/121 (8,3~\%) ; 375–389 : 4/135 (3,0~\%) ; 390–403 : 5/154 (3,2~\%).\par
\noindent Formes relevées (occurrences) : dulcissimae (13), dulcissimo (11), dulcissima (3), carissimae (3), karissimo (1).\par
\noindent{\small Notices (élément attesté sans restitution) : 1, 23, 35, 42, 46, 59, 68, 81, 83, 92, 100, 115, 127, 133, 143, 186, 247, 265, 297, 347, 364, 368, 431, 461.}\par\medskip
\subsection*{plus minus}
\addcontentsline{toc}{subsection}{plus minus}
\noindent Règle (texte normalisé, \textit{u/v} et \textit{i/j} unifiés) : \texttt{\small \textbackslash{}bplus\textbackslash{}s+minus\textbackslash{}b}. Ensemble C : 50 textes avec restitutions, 42 sans restitution ; écart dû aux restitutions : 8.\par
\noindent Par période (attestations sans restitution / textes de la période) : 340–359 : 0/55 (0,0~\%) ; 360–374 : 9/121 (7,4~\%) ; 375–389 : 16/135 (11,9~\%) ; 390–403 : 17/154 (11,0~\%).\par
\noindent Formes relevées (occurrences) : plus minus (49), plus   minus (2).\par
\noindent{\small Notices (élément attesté sans restitution) : 73, 83, 115, 120, 127, 144, 146, 156, 167, 183, 193, 201, 210, 223, 225, 226, 250, 257, 271, 275, 276, 280, 296, 304, 317, 351, 361, 376, 383, 389, 400, 403, 405, 418, 419, 428, 429, 443, 446, 455, 472, 474.}\par\medskip
\subsection*{se vivo / se viva}
\addcontentsline{toc}{subsection}{se vivo / se viva}
\noindent Règle (texte normalisé, \textit{u/v} et \textit{i/j} unifiés) : \texttt{\small \textbackslash{}bse\textbackslash{}s+[ub]i[ub]\textbackslash{}w*}. Ensemble C : 18 textes avec restitutions, 10 sans restitution ; écart dû aux restitutions : 8.\par
\noindent Par période (attestations sans restitution / textes de la période) : 340–359 : 0/55 (0,0~\%) ; 360–374 : 0/121 (0,0~\%) ; 375–389 : 5/135 (3,7~\%) ; 390–403 : 5/154 (3,2~\%).\par
\noindent Formes relevées (occurrences) : se uiui (5), se uiuus (4), se uiuo (3), se uiua (3), se uiuos (1), se   uiuis (1), se uiuum (1).\par
\noindent{\small Notices (élément attesté sans restitution) : 194, 204, 207, 235, 300, 346, 347, 402, 407, 441.}\par\medskip
\subsection*{post consulatum}
\addcontentsline{toc}{subsection}{post consulatum}
\noindent Règle (texte normalisé, \textit{u/v} et \textit{i/j} unifiés) : \texttt{\small \textbackslash{}bpost\textbackslash{}s+con\textbackslash{}w*}. Ensemble C : 23 textes avec restitutions, 15 sans restitution ; écart dû aux restitutions : 8.\par
\noindent Par période (attestations sans restitution / textes de la période) : 340–359 : 4/55 (7,3~\%) ; 360–374 : 2/121 (1,7~\%) ; 375–389 : 8/135 (5,9~\%) ; 390–403 : 1/154 (0,6~\%).\par
\noindent Formes relevées (occurrences) : post consulatum (16), post conssulatum (3), post   consulatu (2), post consulatu (2).\par
\noindent{\small Notices (élément attesté sans restitution) : 15, 16, 17, 34, 102, 109, 179, 183, 184, 185, 186, 261, 275, 301, 350.}\par\medskip

}
